% GENERATED, edit the French unit or the translation table
\mxcorrections{7}\label{C07-CORR}
\textbf{Teacher use.} The answers below concern the tasks not worked in the lesson. First have students state their method; give graded help before the solution. The durations are estimates to be calibrated in class.
\subsection*{Diagnostic and remediation}
\begin{enumerate}
\item $0$. \textit{Action:} Revisit the orthogonality criterion.
\item Yes: $2-2=0$. \textit{Action:} Substitute x and y simultaneously.
\item $x=3$, $y=2$. \textit{Action:} Add, then substitute.
\item $5$. \textit{Action:} Calculate the vector, then its norm.
\item $(x-2)^2-4$. \textit{Action:} Expand to check.
\item The single pair $\pt{0;0}$. \textit{Action:} Use the sum of two non-negative squares.
\end{enumerate}
\subsection*{Investigation activity}
The trials suggest H(2;4), but Pythagoras in PHM proves minimality. The projection is calculated with the normal (2;-1).
\subsection*{Immediate applications}
\begin{enumerate}
\item $3(x-2)+4(y+1)=0$, that is $3x+4y-2=0$.
\item $H\pt{2;1}$ and distance 2; the normal is horizontal.
\item $(x+1)^2+(y-4)^2=0$: the single point $\pt{-1;4}$, not a circle of strictly positive radius.
\item $(x-1)^2=0$, so the only point is $\pt{1;1}$; the line is tangent to the circle.
\end{enumerate}
\subsection*{Exercises: method and results}
\paragraph{Exercise 7.1 -- Constructing a line.} $2(x+1)+5(y-3)=0$, that is $2x+5y-13=0$. Direction $\pt{5;-2}$, so $B\pt{4;1}$ is on the line; check $8+5-13=0$.
\paragraph{Exercise 7.2 -- Calculating a projection.} Normal $\pt{2;-1}$; $\lambda =\dfrac{8-3}{5}=1$. $H=\pt{4;3}-\pt{2;-1}=\pt{2;4}$. $2(2)-4=0$ and $\vect{PH}=\pt{-2;1}$ is normal; distance $\sqrt5$.
\paragraph{Exercise 7.3 -- Recognising a circle.} $(x-3)^2+(y+2)^2=25$; centre $\pt{3;-2}$, radius 5. For $y=-2$, $x=3\pm5$: $\pt{-2;-2}$ and $\pt{8;-2}$.
\paragraph{Exercise 7.4 -- Perpendicular bisector.} Midpoint I(1;2), normal $\vect{AB}=\pt{4;2}$. $4(x-1)+2(y-2)=0$, that is $2x+y-4=0$. I satisfies $2+2-4=0$.
\paragraph{Exercise 7.5 -- The best access.} The projection H(2;4) minimises PH, of length $\sqrt5$ km. For every M on the road, the triangle PHM is right-angled at H and $PM^2=PH^2+HM^2\geq PH^2$. Equality for M=H.
\paragraph{Exercise 7.6 -- Parameter and set.} $(x-1)^2+(y+2)^2=5-k$. If $k<5$, circle with centre $\pt{1;-2}$ and radius $\sqrt{5-k}$; if $k=5$, a point; if $k>5$, empty.
\paragraph{Exercise 7.7 -- A parabola as a locus.} $MF^2=x^2+(y-2)^2$, squared distance to the axis $y^2$. Equality: $x^2-4y+4=0$, that is $y=\dfrac{x^2}{4}+1$. The distances being positive allows the converse.
\paragraph{Exercise 7.8 -- Delayed review.} Centre $\pt{2;0}$, radius 2; $(x-2)^2+y^2=4$. For x=5, $y^2=-5$, impossible. Interpretation: the line is at distance 3 from the centre, greater than the radius.
\paragraph{CQFD, exercise 18 p.~279, questions a, b, c and d.} Possible answers are $\pt{-1;2}$, $\pt{3;-6}$, $\pt{\dfrac14;-1}$ and $\pt{1;0}$. Any nonzero multiple works. For $x=4$, write $x-4=0$.
\paragraph{CQFD, exercise 21 p.~279, questions a and b.} a: $x-(y-2)=0$, that is $x-y+2=0$. b: $4(x+1)+6(y-3)=0$, that is $2x+3y-7=0$. Substituting the coordinates of the points confirms membership.
\subsection*{MCQ: answers and reasons}
\textbf{MCQ 1.}
\begin{enumerate}
\item B. Coefficients of x and y.
\item A. Product with the normal is zero.
\item D. Read the coordinates opposite to the signs.
\item C. Positive root of 9.
\item B. Sum of squares is non-negative.
\end{enumerate}
\textbf{MCQ 2.}
\begin{enumerate}
\item C. Same x-coordinate, y-coordinate zero.
\item A. Absolute vertical gap.
\item D. Nonzero normal.
\item B. Both squares must be zero.
\item C. $y^2=0$.
\end{enumerate}
\subsection*{True or false}
\begin{enumerate}
\item True: through P there is a unique perpendicular to the line; if P belongs to the line, H=P.
\item False: if c=0, a point; if c>0, the empty set.
\item False: substituting a parametrisation of a line into the circle equation gives a quadratic equation with strictly positive leading coefficient, hence at most two roots.
\end{enumerate}
\paragraph{Differentiated teaching.} For a technical block, offer only the first step (domain, factor or reference sample space), then withdraw the help. For an independent student, ask for a second method and the treatment of a limiting case. Orally, allow an initial reformulation in the language the student has mastered, then require a precise mathematical sentence in French.
